\section{Lower-semi-continuity under Rief\/fel quantization}\label{joasa}

We keep the previous setting and inquire now if lower-semi-continuity of the mappings
$t\mapsto\| \P(t)f\|_{\A(t)}$ for all $f\in\A$ implies lower-semi-continuity of the mappings
$t\mapsto\| \PP(t)f\|_{\AA(t)}$ for all $f\in\AA$. We start by noticing that
${\rm Prim}(\A)$ and ${\rm Prim}(\AA)$
are canonically endowed with continuous actions of the group $\Xi$; once again these actions will be denoted by~$\Th$.
By the discussion at the end of Section~\ref{secintra} we are left with proving

\begin{Proposition}\label{segund}
Suppose that $\Q:\C(T)\to \mathcal Z\mathcal M(\A)$ is a covariant $\C(T)$-algebra structure on~$\A$ and that the associated function
$q:{\rm Prim}(\A)\to T$ is open. Then the function $\mathfrak q:{\rm Prim}(\AA)\to T$ associated to
$\mathfrak Q:\C(T)\to \mathcal Z(\AA)$ is also open.
\end{Proposition}

\begin{proof}
We remark f\/irst that $q$ is $\Th$-covariant (Lemma~8.1 in~\cite{Wi}), i.e.\ one has $q\circ\Th_X=q$ for every $X\in\Xi$.
Consequently, if $\mathcal O\subset {\rm Prim}(\A)$ is
an open set, then $\Th_\Xi(\mathcal O):=\{\Th_X(\mathcal K)\mid X\in\Xi, \, \mathcal K\in\mathcal O\}$ will also be an open set
and $q(\mathcal O)=q\left[\Th_\Xi(\mathcal O)\right)]$.
So $q$ will be open if\/f it sends open {\it invariant} subsets of ${\rm Prim}(\A)$
into open subsets of $T$. The same is true for $\mathfrak q:{\rm Prim}(\AA)\rightarrow T$. But the most general
open subset of ${\rm Prim}(\A)$ has the form
\[
\mathcal O_{\mathcal J}:=\{\mathcal K\in {\rm Prim}(\A)\mid \mathcal J\not\subset\mathcal K\}=h(\mathcal J)^c
\]
for some ideal $\mathcal J$ of $\A$, being the complement of the hull  $h(\mathcal J)$ of this ideal.
In addition, $\mathcal O_{\mathcal J}$ is $\Th$-invariant if\/f~$\mathcal J$ is an
invariant ideal. We also recall that Rief\/fel quantization establishes
a~one-to-one correspondence between invariant ideals of $\A$ and invariant ideals of~$\AA$.

So let $\mathcal J$ be an invariant ideal in $\A$ and $\mathfrak J$ its quantization
(an invariant ideal in~$\AA$). We would like to show that
$q\left(\mathcal O_\mathcal J\right)=\mathfrak q\left(\mathcal O_\mathfrak J\right)$;
by the discussion above this would imply that $q$ and $\mathfrak q$ are simultaneously open. Using the fact that
$q(\mathcal K)=t$ if and only if $\mathcal I(t)\subseteq \mathcal K$ and similarly for $\mathfrak q$, one gets
\[
q\left(\mathcal O_\mathcal J\right)=\{t\in T\mid \exists\,\mathcal K\in {\rm Prim}(\A),\ \mathcal J\not\subset\mathcal K,
\ \mathcal I(t)\subset\mathcal K\}
\]
and
\[
\mathfrak q\left(\mathcal O_\mathfrak J\right)=\{t\in T\mid \exists\,\mathfrak K\in {\rm Prim}(\AA),\
\mathfrak J\not\subset\mathfrak K,\ \mathfrak I(t)\subset\mathfrak K\}.
\]
Using the hull application  and the fact that both the hull and the kernel are decreasing, one can write
\[
t\notin q\left(\mathcal O_\mathcal J\right)\ \Longleftrightarrow\ h[\mathcal I(t)]\cap h[\mathcal J]^c=\varnothing
\ \Longleftrightarrow\ h[\mathcal I(t)]\subset h[\mathcal J]\ \Longleftrightarrow\ \mathcal I(t)\supset\mathcal J
\]
and
\[
t\notin \mathfrak q\left(\mathcal O_\mathfrak J\right)\ \Longleftrightarrow\ h[\mathfrak I(t)]\cap h[\mathfrak J]^c=\varnothing
\ \Longleftrightarrow\ h[\mathfrak I(t)]\subset h[\mathfrak J]\ \Longleftrightarrow\ \mathfrak I(t)\supset\mathfrak J.
\]
To f\/inish the proof one only needs to notice that the Rief\/fel quantization of invariant ideals preserves inclusions.
\end{proof}
